% !TEX root = main.tex
\chapter*{Abstrak}

Jerawat (\textit{acne vulgaris}) merupakan penyakit kulit inflamasi kronis yang umum terjadi dan memerlukan penilaian tingkat keparahan (\textit{acne severity grading}) yang akurat guna menentukan intervensi terapeutik yang tepat. Proses diagnosis konvensional secara visual oleh dokter spesialis kulit memiliki keterbatasan berupa kebutuhan waktu yang lama serta rentan terhadap subjektivitas antar-pemeriksa. Pemanfaatan teknologi \textit{deep learning} melalui pendekatan \textit{transfer learning} berbasis \textit{vision foundation model} menawarkan solusi otomatisasi yang efektif untuk mengekstraksi representasi visual kompleks tanpa memerlukan data \textit{training} dalam jumlah masif \textit{from scratch}. Penelitian ini bertujuan untuk melakukan analisis komparatif performa dua skala arsitektur \textit{Vision Transformer} (ViT) berbasis \textit{vision foundation model} \textit{self-supervised} DINOv3 melalui pendekatan \textit{transfer learning}, yang mewakili kategori arsitektur \textit{heavyweight} yaitu DINOv3 ViT-L/16 ($\approx 304$ juta parameter) dan arsitektur \textit{lightweight} yaitu DINOv3 ViT-S/16 ($\approx 21,6$ juta parameter), dalam mengklasifikasikan empat tingkat keparahan jerawat berdasarkan kriteria Hayashi (\textit{Mild}, \textit{Moderate}, \textit{Severe}, dan \textit{Very Severe}; Level 0--3) pada dataset \textit{benchmark} publik \texttt{Layered-Labs/ACNE04} (1.457 citra wajah). Pemilihan keluarga DINOv3 didasarkan pada kemampuannya mempertahankan representasi semantik wajah global sekaligus detail spasial mikrolesi melalui \textit{self-supervised pre-training} pada korpus LVD-1689M serta mekanisme \textit{Gram anchoring}. Eksperimen dirancang secara terstandarisasi menggunakan \textit{preprocessing pipeline} terpadu satu jalur, \textit{stratified data splitting} 70:15:15, dan penanganan \textit{class imbalance}. Kinerja model dievaluasi secara komprehensif menggunakan \textit{confusion matrix} bernotasi angka Level 0--3 serta metrik \textit{Accuracy}, \textit{Macro Precision}, \textit{Macro Recall}, \textit{Macro F1-Score}, dan \textit{Quadratic Weighted Kappa} (QWK), dengan fokus khusus pada sensitivitas kelas sangat parah (Level 3). Hasil dari penelitian ini diharapkan dapat memberikan rekomendasi empiris mengenai \textit{trade-off} antara kapasitas representasi semantik global dan efisiensi komputasi pada arsitektur \textit{vision foundation model} sebagai dasar pengembangan sistem penunjang keputusan dermatologi klinis.

\vspace{0.5 cm}
\begin{flushleft}
    {\textbf{Kata Kunci:} DINOv3, Klasifikasi Jerawat, Transfer Learning, Vision Foundation Model, Vision Transformer.}
\end{flushleft}